\cleardoublepage
\chapter{Analiza wyników badań}
\label{cha:AnalizaBadan}

Celem niniejszego rozdziału jest przedstawienie oraz przeanalizowanie wyników przeprowadzonych eksperymentów. 
Zaimplementowane algorytmy przeanalizowano pod względem czasu sortowania, uzyskanego przyspieszenia i~efektywności, 
wykorzystania zasobów oraz~skalowalności. 
Poza tym uwzględniono również wpływ rozmiaru, typu i~charakterystyki danych wejściowych.

\section{Porównanie czasu sortowania}
\label{sec:CzasSortowania}

Pierwszym i podstawowym kryterium wykorzystanym do oceny wydajności badanych algorytmów był czas wykonania sortowania. 
Analizie poddano warianty sekwencyjne oraz równoległe algorytmów Quick Sort, Merge Sort, 
Bucket Sort i~Sample Sort dla zbiorów danych o rozmiarach 1~000, 10~000, 100~000 oraz 1~000~000 elementów.

W pierwszej kolejności przeanalizowano wpływ rozmiaru danych na czas wykonania poszczególnych algorytmów. 
Jako punkt odniesienia przyjęto zbiór danych \texttt{random\_int}, zawierający losowo wygenerowane liczby całkowite. 
Pozwala to na porównanie badanych algorytmów w jednakowych warunkach, 
bez uwzględniania na tym etapie wpływu stopnia uporządkowania danych lub występowania duplikatów.

\begin{figure}[htbp]
    \centering
    \includegraphics[width=0.90\textwidth]
    {images/chapter6/execution_time/execution_time_vs_data_size_random_int.png}
    \caption{Średni czas wykonania wariantów sekwencyjnych w zależności od rozmiaru zbioru \texttt{random\_int}.}
    \label{fig:time-data-size-sequential}
    
    {\small Źródło: opracowanie własne.}
\end{figure}

Dla wszystkich badanych algorytmów zwiększenie rozmiaru zbioru skutkowało wzrostem czasu wykonania, 
przy czym różnice dla różnych algorytmów szczególnie stawały się widoczne dla największego badanego zbioru.

Zestaw danych obejmujący 1~000~000 losowych liczb całkowitych uzyskał najkrótszy średni czas wykonania dla algorytmu Bucket Sort, 
osiągając około 2,04~s. Nieco dłuższego czasu wymagał Quick Sort z czasem około 2,25~s. 
Merge Sort wymagał około 3,41~s, natomiast Sample Sort osiągnął najdłuższy czas spośród porównywanych algorytmów, równy 4,95~s.
Wyniki wskazują, że różnice pomiędzy badanymi implementacjami zwiększają się wraz ze wzrostem rozmiaru przetwarzanego zbioru.

W celu bezpośredniego porównania wariantów sekwencyjnych i~równoległych w tabeli~\ref{tab:best-execution-time} 
zestawiono czasy uzyskane dla 1~000~000 losowych liczb całkowitych. 
Dla każdego algorytmu przedstawiono czas wariantu sekwencyjnego oraz najlepszy czas uzyskany spośród badanych konfiguracji równoległych.

\input{tables/chapter6/best_execution_time}

Zestawienie pokazuje, że dla każdego z badanych algorytmów możliwe było uzyskanie czasu krótszego niż w wariancie sekwencyjnym, 
niemniej najlepsze wyniki osiągęto przy różnej liczbie jednostek wykonawczych. 
Największe przyspieszenie w stosunku do wykonania sekwencyjnego odnotowano dla Sample Sort i~wyniosło ono około 2,23. 
Nie oznaczało to jednak najkrótszego bezwzględnego czasu sortowania. 
Najlepszy wynik czasowy uzyskał Bucket Sort, dla którego przy wykorzystaniu 16 jednostek wykonawczych wyniósł on około 1,09~s.

Uzyskane wyniki wskazują, że zastosowanie wariantu równoległego może prowadzić do skrócenia czasu sortowania,
jednak wielkość uzyskanej poprawy jest zależna od badanego algorytmu oraz zastosowanej konfiguracji równoległości.

\section{Wpływ liczby jednostek wykonawczych na wydajność}
\label{sec:LiczbaRdzeni}

Kolejnym analizowanym aspektem był wpływ liczby wykorzystywanych jednostek
wykonawczych na czas działania równoległych wariantów badanych algorytmów.
Analizę przeprowadzono dla zbioru \texttt{random\_int} zawierającego 1~000~000 elementów, 
wykorzystując kolejno 2, 4, 8, 16 oraz 32 jednostki wykonawcze. 
Pozwoliło to przeanalizować, jak zmiana liczby jednostek wykonawczych wpływa na czas sortowania.

W celu porównania wpływu liczby jednostek wykonawczych na wydajność poszczególnych algorytmów 
zestawiono wyniki uzyskane dla wszystkich badanych konfiguracji.

\input{tables/chapter6/scalability_time}

Uzyskane wyniki wskazują wyraźne różnice pomiędzy badanymi algorytmami.
Wzrost liczby jednostek wykonawczych nie zawsze przekładał się na poprawę uzyskiwanego czasu sortowania, a najkrótsze
czasy osiągano przy różnej liczbie wykorzystywanych jednostek.

W przypadku Quick Sort najkrótszy średni czas wykonania uzyskano już przy wykorzystaniu 2 jednostek wykonawczych i~wyniósł on około 1,83~s.
Zwiększenie ich liczby do 4 skutkowało wydłużeniem czasu do około 2,22~s, natomiast dla 8 jednostek osiągnął on około 2,56~s. 
Dalsze zwiększanie liczby jednostek do 16 i 32 nie powodowało już wyraźnych różnic w czasie wykonania, 
który utrzymywał się na poziomie około 2,58~s. 
Wyniki wskazują, że dla badanej implementacji Quick Sort zwiększenie liczby jednostek wykonawczych 
powyżej 2 nie przynosiło dalszych korzyści pod względem czasu wykonania.

Dla Merge Sort zwiększenie liczby jednostek wykonawczych z 2 do 4 pozwoliło skrócić czas z około 2,34~s do 2,16~s. 
Dalsze zwiększenie ich liczby do 8 nie przyniosło dalszej poprawy, a czas wyniósł około 2,32~s.
Dla 16 oraz 32 jednostek wykonawczych odnotowano wyraźne wydłużenie czasu wykonania, który wyniósł odpowiednio około 3,78~s oraz 3,85~s.
Najkorzystniejszy wynik dla Merge Sort osiągnięto więc przy 4 jednostkach wykonawczych.

Odmienną tendencję zaobserwowano dla Bucket Sort. 
Zwiększanie liczby jednostek wykonawczych od 2 do 16 prowadziło do systematycznego skracania czasu sortowania. 
Średni czas skrócił się z około 1,76~s dla 2 jednostek do 1,52~s dla 4, 1,19~s dla 8 oraz 1,09~s dla 16 jednostek wykonawczych.
Dopiero zwiększenie ich liczby do 32 skutkowało wydłużeniem czasu do około 1,25~s. 
Spośród badanych algorytmów Bucket Sort wykazywał poprawę czasu wykonania przy zwiększaniu liczby jednostek wykonawczych aż do 16 jednostek.

Sample Sort również początkowo charakteryzował się poprawą wydajności.
Średni czas wykonania zmniejszył się z około 3,93~s dla 2 jednostek do 2,76~s dla 4 oraz 2,22~s dla 8 jednostek wykonawczych. 
Po zwiększeniu ich liczby do 16 nastąpiło jednak wyraźne wydłużenie czasu wykonania do około 4,96~s. 
Przy 32 jednostkach wartość ta uległa jedynie niewielkiej zmianie osiągając około 4,99~s. 
Najlepszy wynik uzyskano zatem przy 8 jednostkach wykonawczych.

Zmiany czasu wykonania poszczególnych algorytmów wraz ze zwiększaniem liczby 
jednostek wykonawczych przedstawiono na rysunku~\ref{fig:time-cores-random-int}.

\begin{figure}[htbp]
    \centering
    \includegraphics[width=0.90\textwidth]
    {images/chapter6/execution_time/execution_time_vs_cores_random_int.png}
    \caption{Średni czas wykonania wariantów równoległych dla zbioru
    \texttt{random\_int} o rozmiarze 1~000~000 elementów w zależności od liczby jednostek wykonawczych.}
    \label{fig:time-cores-random-int}

    {\small Źródło: opracowanie własne.}
\end{figure}

Uzyskane wyniki wskazują, że najkorzystniejsza liczba jednostek wykonawczych różniła się w zależności od badanego algorytmu.
Quick Sort osiągnął najlepszy wynik przy 2 jednostkach, Merge Sort przy 4, 
Sample Sort przy 8, natomiast Bucket Sort przy 16 jednostkach wykonawczych.
Zwiększanie ich liczby powyżej tych wartości nie prowadziło do dalszej poprawy, 
a dla niektórych algorytmów wiązało się ze wzrostem czasu wykonania.

W celu określenia, czy podobne zależności występują również dla innego typu danych, przeanalizowano wyniki dla zbioru \texttt{random\_float}.
Dla Merge Sort, Bucket Sort oraz Sample Sort najkrótszy czas wykonania uzyskano przy tej samej liczbie jednostek wykonawczych 
co dla zbioru \texttt{random\_int}, odpowiednio przy 4, 16 oraz 8 jednostkach. 
Wyjątek stanowił Quick Sort, dla którego najkrótszy czas przy danych zmiennoprzecinkowych uzyskano przy 4 jednostkach wykonawczych zamiast 2.

Podobną tendencję zaobserwowano szczególnie dla Bucket Sort oraz Sample Sort.
Dla obu typów danych Bucket Sort skracał czas sortowania do 16 jednostek,
natomiast Sample Sort najkrótszy czas uzyskiwał przy 8 jednostkach. 
Zwiększenie ich liczby do 16 i~32 prowadziło już do wyraźnego wydłużenia czasu wykonania.
Powtarzalność tych zależności dla dwóch typów danych sugeruje, że obserwowane zmiany wydajności mogły wynikać 
przede wszystkim ze sposobu realizacji równoległości badanych algorytmów, a nie wyłącznie z typu przetwarzanych wartości.

Uzyskane wyniki pokazują, że zwiększenie liczby wykorzystywanych jednostek nie prowadziło do proporcjonalnego wzrostu wydajności.
Znaczenie ma również sposób zrównoleglenia konkretnego algorytmu oraz narzut związany z organizacją obliczeń równoległych. 
Wpływ tych czynników można ocenić również na podstawie wykorzystania zasobów oraz uzyskanego przyspieszenia i~efektywności.

\section{Analiza zużycia pamięci RAM i obciążenia CPU}
\label{sec:RamCpu}

Kolejnym analizowanym aspektem było wykorzystanie zasobów sprzętowych podczas działania równoległych wariantów badanych algorytmów. 
Uwzględniono średnie obciążenie procesora oraz średnie i~maksymalne wykorzystanie pamięci RAM. 

W celu bezpośredniego porównania wykorzystania zasobów przez badane algorytmy
przyjęto wspólną konfigurację wykorzystującą 8 jednostek wykonawczych,
stanowiącą środkową wartość spośród badanych konfiguracji.

W tabeli~\ref{tab:resources} przedstawiono średnie obciążenie CPU oraz wykorzystanie pamięci RAM dla 
zbioru \texttt{random\_int} zawierającego 1~000~000 elementów.

\input{tables/chapter6/resources}

\subsection*{Obciążenie procesora CPU}

Wartości obciążenia CPU mogą przekraczać 100\%, ponieważ pomiar uwzględnia zarówno proces główny, 
jak i~procesy potomne wykorzystywane podczas wykonywania algorytmu równoległego.
W zastosowanej metodzie pomiarowej 100\% oznacza pełne obciążenie jednego procesora logicznego, 
natomiast wyższe wartości wskazują na równoczesne wykorzystanie kilku jednostek wykonawczych.

Przy wykorzystaniu 8 jednostek wykonawczych najwyższe średnie obciążenie CPU odnotowano dla Sample Sort i~wyniosło ono około 431,51\%.
Drugą najwyższą wartość uzyskał Bucket Sort z wynikiem około 365,50\%. 
Merge Sort oraz Quick Sort charakteryzowały się wyraźnie mniejszym obciążeniem procesora, osiągając odpowiednio około 182,60\% i~149,62\%.

Zmiany obciążenia procesora wraz ze wzrostem liczby jednostek wykonawczych przedstawiono na rysunku~\ref{fig:cpu-cores-random-int}.
Dodatkowo zaznaczono teoretyczne maksymalne obciążenie CPU, odpowiadające 100\% dla każdej wykorzystywanej jednostki wykonawczej.

\begin{figure}[htbp]
    \centering
    \includegraphics[width=0.95\textwidth]
    {images/chapter6/cpu/cpu_vs_cores_comparison_random_int.png}
    \caption{Średnie wykorzystanie CPU dla zbioru \texttt{random\_int}
    o rozmiarze 1~000~000 elementów w zależności od liczby jednostek
    wykonawczych.}
    \label{fig:cpu-cores-random-int}

    {\small Źródło: opracowanie własne.}
\end{figure}

Dla Bucket Sort zwiększanie liczby jednostek wykonawczych skutkowało stopniowym wzrostem obciążenia procesora. 
Średnie wykorzystanie CPU wzrosło z około 165,35\% dla 2 jednostek do 231,69\% dla 4,
365,50\% dla 8 oraz 547,26\% dla 16 jednostek wykonawczych.
Największe średnie obciążenie CPU odnotowano przy 32 jednostkach wykonawczych i~wyniosło ono około 768,07\%.
Wyniki pokazują, że wraz ze zwiększaniem liczby jednostek wykonawczych Bucket Sort generował coraz większe łączne obciążenie procesora.

Wzrost obciążenia CPU nie prowadził jednak w każdym przypadku do dalszego skracania czasu sortowania. 
Dla Bucket Sort najlepszy czas uzyskano przy 16 jednostkach wykonawczych, natomiast zwiększenie ich liczby do 32 spowodowało 
wydłużenie średniego czasu z około 1,09~s do 1,25~s, pomimo wzrostu średniego wykorzystania procesora z około 547\% do 768\%. 
Wskazuje to, że większe wykorzystanie CPU nie zawsze wiąże się z poprawą czasu wykonania.

Odmienną tendencję zaobserwowano dla Merge Sort. 
Wykorzystanie CPU wzrosło z około 184,96\% dla 2 jednostek wykonawczych do około
247,28\% dla 4 jednostek, po czym przy 8 jednostkach spadło do około 182,60\%. 
Dla 16 oraz 32 jednostek wykonawczych średnie obciążenie procesora osiągało już tylko około 98\%. 
Zmniejszeniu obciążenia CPU towarzyszył jednocześnie wzrost czasu sortowania. 
Najkrótszy średni czas wykonania Merge Sort odnotowano dla 4 jednostek wykonawczych, 
natomiast przy 16 i~32 jednostkach wzrósł on do około 3,8~s.

Podobną zależność zaobserwowano dla Quick Sort. 
Wyższe obciążenie procesora występowało przy mniejszej liczbie jednostek wykonawczych. 
Dla 2 jednostek wynosiło ono około 170,36\%, natomiast dla 4 osiągnęło około 179,53\%.
Przy 8 jednostkach wartość zmniejszyła się do około 149,62\%, a dla 16 i~32 wynosiła już około 98\%.
Wraz ze zwiększaniem liczby jednostek wykonawczych nie obserwowano odpowiadającego mu wzrostu obciążenia procesora.

Wyraźne zmiany wystąpiły również dla Sample Sort. 
Średnie obciążenie procesora wzrosło z około 191,57\% dla 2 jednostek
wykonawczych do 309,72\% dla 4, a~najwyższy poziom około 431,51\% odnotowano przy 8 jednostkach. 
Dalsze zwiększenie ich liczby do 16 i~32 skutkowało zmniejszeniem wykorzystania CPU do około 99\%.

Porównanie wyników wskazuje zatem na wyraźne różnice w sposobie wykorzystania dostępnej równoległości przez poszczególne algorytmy.
Bucket Sort jako jedyny zwiększał średnie obciążenie CPU w całym badanym zakresie od 2 do 32 jednostek wykonawczych.

Kolejnym czynnikiem wpływającym na wykorzystanie procesora był rozmiar przetwarzanego zbioru. 
Zależność średniego obciążenia CPU od liczby sortowanych elementów dla 32 jednostek wykonawczych 
przedstawiono na rysunku~\ref{fig:cpu-size-random-int}.

\begin{figure}[htbp]
    \centering
    \includegraphics[width=0.95\textwidth]
    {images/chapter6/cpu/cpu_vs_data_size_random_int.png}
    \caption{Średnie wykorzystanie CPU w zależności od rozmiaru zbioru
    \texttt{random\_int} dla 32 jednostek wykonawczych.}
    \label{fig:cpu-size-random-int}

    {\small Źródło: opracowanie własne.}
\end{figure}

Dla najmniejszych zbiorów danych w części pomiarów odnotowano bardzo niskie,
a w niektórych przypadkach również zerowe wartości średniego wykorzystania CPU.
Wyniki te należy interpretować z uwzględnieniem przyjętego sposobu próbkowania obciążenia procesora.
Ze względu na krótki czas wykonania sortowania dla najmniejszych zbiorów pomiar mógł zakończyć się przed zarejestrowaniem
wystarczającej liczby próbek. Wartość równa 0\% nie oznacza zatem braku wykorzystania procesora, lecz wynika 
z ograniczenia zastosowanej metody pomiarowej.
Dla większych zbiorów danych rejestrowane wartości wykorzystania procesora były wyższe. 
Szczególnie widoczne było to dla Bucket Sort, dla którego przy 32 jednostkach wykonawczych średnie wykorzystanie CPU wzrosło
z około 65,63\% dla 100~000 elementów do około 768,07\% dla 1~000~000 elementów.

Dla Merge Sort, Quick Sort oraz Sample Sort przy 32 jednostkach wykonawczych i~1~000~000 elementów 
średnie wykorzystanie procesora wynosiło natomiast około 97--99\%. 
Oznacza to, że w tych konfiguracjach zwiększona liczba jednostek wykonawczych nie przełożyła się na wyższe obciążenie CPU. 
Zaobserwowane różnice pomiędzy algorytmami wskazują na znaczenie sposobu organizacji obliczeń równoległych.

Zbliżony przebieg zaobserwowano również dla zbioru \texttt{random\_float}. 
Bucket Sort zwiększał wykorzystanie CPU wraz ze wzrostem liczby jednostek wykonawczych, 
natomiast dla Sample Sort najwyższe obciążenie procesora wystąpiło przy 8 jednostkach.
Wskazuje to, że zaobserwowane zależności nie były charakterystyczne wyłącznie dla danych typu całkowitego. 

\subsection*{Wykorzystanie pamięci RAM}



\section{Skalowalność algorytmów}
\label{sec:AnalizaSkalowalnosci}

Ocena, jak dobrze algorytmy skalują się wraz ze wzrostem liczby rdzeni i rozmiarem danych

\section{Wpływ typu i charakterystyki danych wejściowych}
\label{sec:CharakterystykaDanych}

Porównanie wyników dla tych samych algorytmów, ale różnych typów danych wejściowych


Omówienie wyników dla danych losowych itp. Wskazanie, który algorytm radził sobie lepiej przy danym rodzaju danych